Software systems spanning file parsers, multimedia engines, and asynchronous network servers present diverse input domain barriers for coverage-guided fuzzers.
In the \drogon web framework, an uncaught filesystem exception inside a method marked \texttt{noexcept} triggered process aborts on oversized upload filenames.
We present \tool (Quick Transfer Fuzzer), an adaptive fuzzing architecture combining structure-aware mutation, an independent multi-armed bandit scheduler (Discounted-UCB), and headless in-memory execution shims.
We designed a base structural mutator for \fastgltf, then extended the architecture to multimedia rendering in \raylib via software rasterizer shims, and to asynchronous network protocols in \drogon via in-memory socket mocks.
Across 59 target binaries, \tool executed 807{,}520{,}128 inputs, reaching up to 53{,}732 executions per second on network parsers without operating system socket bottlenecks.
\tool uncovered 16 security defects, with six upstream-confirmed vulnerabilities resulting in two merged pull requests and four maintainer-committed patches.
Our results show that domain-specialized in-memory shims and non-stationary bandit adaptation enable effective fuzzing across heterogeneous software domains.
